\section{Julius Evola on Esoteric Catholicism}

\begin{quotex}
Today with difficulty, if not exceptionally in some close to dangerous existential crises, the potentiality of Christianity at its beginnings as that ``tragic doctrine of salvation'' can be re-actualized. The problem is not set and we even say without reticence that if anyone who has known, for some time, nothing other than the vainest constructions of philosophy and the secular plebeian university culture of today or the contaminations of the various contemporary individualisms, aestheticisms, and romanticisms, would ``convert'' to Catholicism and would experience truly the faith at least, with a total commitment and possibly in a ``sacrificial'' sense, that would signify not an abdication but rather, in spite of everything, a progress.

\flright{From \emph{Maschera e Volto dello Spiritualismo Contemporaneo}}
\end{quotex}

\flrightit{Posted on 2011-03-31 by Aeneas}

\begin{center}* * *\end{center}

\begin{footnotesize}
\begin{sffamily}
\texttt{Cologero on 2012-12-21 at 09:20 said:}

Alex, we here at Gornahoor have gone through a lot of trouble to make Evola's work available, including much of his correspondence. At this point, there are several hundred thousand words of text, but there is a search feature available plus google indexes everything here.

Since you seem to be a recent visitor, I'll make you aware of this text by Julius Evola:
\textit{Esoteric Catholicism and Integral Traditionalism}\footnote{\url{http://www.gornahoor.net/library/EsotericCatholicism.pdf}}. So obviously your grasp of his work has limitations. This is not a criticism, since we all have limiatations; keep that in mind whenever and wherever you read anything.

We have elucidated several hermeneutic principles. First is that, regarding Guenon and Evola, Guenon is the Master and Evola the pupil. This is not just our position, Evola himself admits it. In their correspondence over decades, Evola went to Guenon for advice, but never the reverse. A careful reading of our texts will reveal inconsistencies in Evola over the years. Evola seemed to have, as Guenon points out, certain prejudices that colored his opinions, as well as ambitions to be academically acceptable, something Guenon did not care about. Also, Evola was never initiated into anything, so we must take his ipse dixit pronouncements with a grain of salt.

So let's start with Guenon. He addresses the Christian tradition in several of his works and understands its esoteric origins. Moreover, he points out that such works as Dante's Divine Comedy and the Romance of the Rose are esoteric texts. So anyone with the proper understanding can certainly see that there definitely is an ``esoteric dimension'' in Catholicism, at least as it was understood in the Middle Ages, which is really all we are interested in here.

Another hermeneutic principle is to look at whether a teaching is consistent with the metaphysics of Tradition. We see that in the Middle Ages, there was an understanding of the different levels of the soul, that there is an intellect that is superior to the rational mind, that there is a cosmic order, that the persons in society fall generally into a caste-like division of functions, and so on. We see this all in the spiritual tradition of the Middle Ages. Not even Evola denies that the European Middle Ages was a Traditional society. Obviously, this was not the result of an accident, but indicates the existence of an esoteric tradition that manifested in what we see as its outward forms.

Ultimately, how does it help you if someone tells you there was an esoteric society in Prague in 1647 or some such thing? This is all mechanical thinking and can only result in the simulacra of authentic esoteric understanding. Your starting point must always be self-knowledge and self-understanding, to see how your mind works, where your thoughts come from, how desires arise, and so on. As your own consciousness becomes more ordered, it will be easier to understand esoteric texts.

\hfill

\texttt{Jason-Adam on 2012-12-22 at 01:47 said:}

I apologise if I intrude but several of Alex's points are in error :

1. The Holy Grail is not of pagan origin but is an esoteric representation of the eucharist. I recommend strongly reading AE Waite's work of the Grail to understand its role as a symbol of inner Christianity.

2. Guenon did admit the existence of initiation in Christianity but that it was no longer open to all as the organisations in question were withdrawn from the world. With time that could change.

3. Evola was indeed quite inconsistent because at the end of his life, in Orientations, he admitted that the Church held to the positions of the Syllabus of Errors, then he would support the Church. In other words Evola was of the same opinion as the late Archbishop Lefebrve.

4. Alex, I believe your development would be aided by reading the works of Jean Borella...

\hfill

\texttt{Mihai on 2012-12-22 at 09:56 said:}

The Grail has its antecedents in ancient Celtic legends and traditions and it was considered a pre-figuration of the Eucharist, thus transposed in Christian terminology.

Also, Waite is an author that can hardly be taken seriously, given his connections with ``organisations'' such as The Golden Dawn.

Guenon actually reviewed that very book in an article which is part of \textit{Symbols of the sacred science}.

\hfill

\texttt{Cologero on 2012-12-22 at 17:47 said:}

Alex: Guenon has written extensively about initiatic organizations in the West, which he claimed existed up until a few centuries ago. He also made a rather surprising revelation in one of his letters to Evola. Guenon also explained his conversion to Islam, which is probably not what you think. All that is documented here, and in his books, so perhaps you can comment on specifics.

You have a misunderstanding of initiation or tradition, by confusing the Tradition with the various outward forms it may take. What is your purpose in that if not polemical? If you want to strut your stuff, then explain why the Divine Comedy is not indicative of initiatory knowledge? Guenon certainly thought it was, and if he was wrong about that, how can we trust him on anything else?

Evola, unfortunately, mixes up non-traditional ideas with his own philosophical ``system''. That is what makes him appealing, and we add, comprehensible, to certain kinds of minds. He is good on the criticism of the modern world, but unreliable on anything of a positive nature.

\hfill

\texttt{Logres on 2012-12-22 at 23:07 said:}

Alex said ``That doesn't mean they were contained inside the main Christian religion... ''

I don't understand this point of view -- it sounds as if you are saying ``unless so-and-so exoteric religion has organized within itself, as an institution (albeit a secret one, or a ``hidden one''), which is sort of an acknowledged `state secret', then that religion cannot be said to `contain' an esoteric tradition''.

Those seem like very loaded terms -- granted, no one is arguing that the exoteric Christian religion did everything it could to foster esotericism in the West -- if it had done this, it would either

a) be esoteric itself

b) the exclusive product of said esoteric society

I am not sure if any religion would fit category b, and choice a seems to undercut the idea of ``religion'' to begin with.

Instead, Christianity seems to have been a hybrid or amalgamation of sorts, a ``Platonism for the masses'' or (as Borella put it) a sort of esotericism made into exoteric religion. Perhaps one can disagree with the ``gambit'' of doing this, but I am not sure what the point would be -- every possibility will have to manifest, and this possibility was certainly latent in any attempt to inter-relate esoteric mystery with exoteric dogma.

\hfill

\texttt{Jason-Adam on 2012-12-22 at 23:27 said:}

As a way of relation, it must be said that the esoteric-exoteric conflict which exists in Christianity also exists in other traditional belief systems -- in Islam many exoteric Sunnis and Shia believe Sufism to be heretical and Sufis have been persecuted from time to time. In China and Korea exoteric Confucians at time persecuted esoteric Taoists and Buddhists. Evolians are holding Christianity to an imaginary high standard no religion in the world is up to par on. Ultimately we are dealing with the prejudice of a Dadaist and his heavy metal fan followers who despise their ancestries and cover up their nihilism in pseudo-academic effete jargon.

Now to return to the nature of Christianity, it is indeed Platonism for the masses as a philosophy professor of mine also put it, Schuon also agreed on the notion of esoteric made exoteric. More to the point, Boris Mouravieff shows that Christianity is not simply another form of the tradition but it is the completion and fulfillment of all traditions. Christianity is the universal religion to unite all men under God's kingdom.

One word of advice to students I should add -- instead of focusing solely on politics as Evolians tend to do, one should study the art of a particular time and place to know the inner feeling and dynamic of said society. The Christian Mediaeval West had art based on the theme of transcendence, whereas the Classical art preceeding and the Renaissance art following it were human-based art styles. Art history matches prefectly with Guenon's interpretation of Christianity as restorer of Western tradition from Greco-Roman degeneracy and the decline of the Church is the decline of the west... .

\hfill

\texttt{Mihai on 2012-12-23 at 06:09 said:}

@Alex:

Try to look a little into what the Christian East contains.

Try Dionysius the Areopagite, Maxim the Confessor, Gregory Palamas, John of the Ladder as well as many other saints whose teachings have been collected in the Philokalia. The Hesychastic Way, the way of the Heart, is the very core of Christianity and represents an integral Christian esoterism, without it being mixed with non/pre-Christian elements, as was the case with Christian Hermeticism and the Grail.

I agree that Catholicism has tended, with each passing century, to cultivate a more and more rationalist approach to religion, which lead to the numerous problems that have sprung in the West, beginning with the very idea behind modernity itself.

Also, you have to understand that Evola is no authority when discussing the Christian tradition. Some of the arguments he brings in that chapter posted by Cologero are of the kind that only a rationalist academic critic would make.

Evola's approach to spirituality, also has some flaws, some only in particulars, but a few in essential matters, which cannot be passed over.

One such example is his view of spiritual as something ``impersonal''.

\hfill

\texttt{Logres on 2012-12-23 at 15:01 said:}

Alec, those are good recommendations by Mihai -- everyone's been argued with on this Forum -- it's part of who we are -- iron sharpening iron. You might be interested in the Hesychasm controversy that erupted around Barlaam, who denied ``navel-gazing'' : the Eastern Church strongly defended not only the ideal, but the specific practices.
\url{http://en.wikipedia.org/wiki/Barlaam_of_Seminara}

Cologero has done a lot of work digging up obscure Christian teachings and examining them in light of Tradition -- the consensus on this Forum is against you, so it's a rough, uphill row to how, because it's not a religious, but an intellectual consensus, between those of widely varied practice and background.

Cologero is right to emphasize a certain purity of will and intuition here -- you find what you seek. What needs to happen, as soon as possible, is that Traditionalists begin reweaving the garment of God, which has been ripped asunder. It exposes (to the casual eye) the project to the accusation that we are cherry-picking for a special thesis, but this is something deeper than a will-based and relativistic reshaping of history ``as we would like it to be'' -- it is the Hermetic method of seeing below the surface, and discerning the Logos at work. The Church has a lot to answer for, and judgement starts at the household of God; but we have the critical evidence and the criticism, which only sets the task before us: to see deeper into the mysteries. As you do so, you will no doubt come upon ``things'' others are not aware of -- test the spirits, and keep going. But don't take a starting point as an ending point, and shut the Book. Keep going.

\hfill

\texttt{Logres on 2012-12-23 at 15:03 said:}

Also, as regards Evola, one learns (or should learn) more from the errors of genius, than the ``correct thinking'' of the mediocre. That is the approach one should have with Evola.

\hfill

\texttt{Senko on 2012-12-23 at 21:38 said:}

Alex, Charbonneau-Lassay, who was a really orthodox Catholic was the depositary of the Order of the Paraclite, that was a ``mainstream'' Catholic initiatic order. By the way, in this issue I have to agree with Reyor and Borella, the very Sacraments are meant to be initiatic, so, for those qualified and aware of their nature they are initiatic, for the rest, they appear as exoteric, though the initiatic influences are still there, even if wasted or not ``used''.

\hfill

\texttt{Senko on 2012-12-23 at 21:41 said:}

Mihai,

You are right pointing to the Orthodox East, but if you let me, I'll add that all this can be used in the West too. My parrish priest (I'm Catholic) taught me Prayer of the Heart, and I used to discuss with him so Orthodox and Hermetic authors (even Guenón a few times)

\hfill

\texttt{Logres on 2012-12-23 at 22:16 said:}

Alex: I am not so sure you just asked a question, you also actually made the assertion:

``I certainly agree the European Middle Ages was a traditional society, but Christianity in itself never had an initiatic system... ''

\hfill

\texttt{Mihai on 2012-12-24 at 02:12 said:}

You are certainly lucky to have found a parish priest who is open towards things like Hermetism and esoteric matters in general.

Nowadays, Christian priests (both in the East and in the West) who have an integral knowledge of the very heart of Christian teachings are hard to find, but not impossible.

\hfill

\texttt{Jason-Adam on 2012-12-27 at 00:33 said:}

1. When dealing with matters of truth or falsehood emotions play no role in the decision, Julius Evola was a seeker after truth and a lover of order, as such we are all united in the same endeavour, thus I respect Baron Evola as a man who offered sometime profound insights. However he sadly was infected with the spirit of rebellion which created the modern world he despised and he never quite shook off his Dadaist past, his hatred of Christianity being part of said rebellion. However, at various times in his life, he seems quite close to the truth. I've learned over the years to sense a writer's inner feelings that he tries to disguise in his words and I can tell that for all his Stoic front he put on, Evola was hollow and looking for what he could not find. If only he had lost a little pride, as God Himself suggested with the wound in 195, he could have found liberation in Catholicism.

The respect I have for a tragic hero such as Evola does not extend to those degenerate pariahs who use his name to justify their antisocial and destructive behaviour. It must be pointed out, as I am surprised Cologero hasn't yet, that Guenon clearly said in La crise du monde moderne that a dead tradition can not be revived once an organisation has no successors to its spiritual power -- we can all agree that paganism has no valid chain of initiation that go back to pre-Christian times -- and as such Evola did not advocate restoring something he knew could not be done but instead he believed in abstract principles with the hope that revelation would come at a future date.

And one more note -- in the piece Cologero translated above Evola actually said that on the vulgar level it is correct to defend Catholicism from other spiritualities in the West for foreign beliefes divorced from their places of origin are prey to becoming agents of the countertradition, as we can see today in various perversions of Brahmanism, Buddhism etc etc... .

\hfill

\texttt{Cologero on 2012-12-27 at 06:38 said:}

Very well put, Jason-Adam. It is not up to me to say everything, but anyone with an objective mind can put things together as you have. On the positive side, Evola has been instrumental in leading many to see the shortcomings of the modern world. That, of course, includes the existing spiritual traditions which seem far from what they used to be. I believe that I have documented sufficiently how the spiritual worldview of 1000 years ago differs significantly from what is commonly held today. Perhaps you can understand why Evola was so reluctant to associate himself with its contemporary forms.

Unfortunately, that left him in a precarious position, a ``traditionalist'' with no Tradition. He knew paganism could not be restored; that is why he was so harsh in his criticism of neo-paganism. From his correspondence with Guenon, we can see that even into his 50s, he was still searching for an initiation or tradition, even through marginal figures. That is sad.

His attitude toward Dadaism is quite strange. Clearly, it was a revolutionary movement; we can understand why a young man would have been attracted to it, but would have dropped it after finding Tradition. He seems to have done so, since he associated Dada, and its Romanian founder, with ``cultural Marxism'', as we showed in an essay written in the 1930s that we translated. Yet in his spiritual autobiography, he retracts himself and goes on to praise Tristan Tzara. Then there is that rather sad interview with him floating around on youtube where he shows great joy in his association with Dada. To shock the bourgeoisie is a left wing sport.

Evola often wrote about the inner spiritual combat. It should not be necessary to point out that there is an abyss between writing about it and actually engaging in it. At this point, it is no longer sufficient to simply quote Evola as an ``inspiration''; our inspiration needs to come from a higher source.

\hfill

\texttt{Jason-Adam on 2012-12-27 at 17:49 said:}

One blessing in disguise I think Vatican II has brought is that in the fight against it, traditionalist Catholics have been inspired to rebirth the Mediaeval spirit in themselves and not remain passive as too many Catholics had become by the 1930s-1940s. While the circle of traditional Catholicism is quite tiny, we do not know yet if they may amount to something more in the future... ... ... ... ..

Evola I think had 3 phases in his life -- the initial revolutionary period in the 20s, his shift to conservatism in the 30s and 40s, and then his ``anarchist/active nihilist'' post-war return to his youth in the 50s and 60s. Even yet though, as late as 1974 he was writing that he hoped the Church could hold Spain together after Franco's death... ... \end{sffamily}
\end{footnotesize}

